\subsection{\texorpdfstring{Application: relations between the spaces \(\mathcal{L}_{F}^{p}\) ( \(1 \leqslant p \leqslant +\infty\) )}{Application: relations between the spaces L sub F to the p ( 1 <= p <= +infinity )}}

PROPOSITION 4. --- Let f be a measurable function with values in a Banach space F; the set I of numbers p, such that \(1 \leqslant p \leqslant +\infty\) and \(\mathrm{N}_{p}(\mathbf{f})\) is finite, is either empty or is an interval of \(\overline{R}\) . If I is nonempty, the restriction to I of the mapping \(p \mapsto \mathrm{N}_{p}(\mathbf{f})\) is continuous; moreover, if f is not negligible, \(\log \mathrm{N}_{p}(\mathbf{f})\) is a convex function of 1/p on \(\overline{I}\) .

We already know (Ch. I, No.~3, Prop. 5) that the set J of finite numbers \(p \geqslant 1\) such that \(\mathrm{N}_{p}(\mathbf{f}) < +\infty\) is either empty or is an interval, and that \(\log \mathrm{N}_{p}(\mathbf{f})\) is a convex function of 1/p on J (when f is not negligible); this of course implies the continuity of \(p \mapsto \mathrm{N}_{p}(\mathbf{f})\) on J.

If J is empty then either \(I = \varnothing\) or \(I = \{+\infty\}\) , and the proposition is obvious in this case; assume henceforth that J is nonempty. The proposition is also obvious if f is negligible; assume henceforth that f is not negligible. If \(s \in J\) then, for every finite number p \textgreater{} s, \(|f|^{p} = |f|^{s}|f|^{p-s}\) , and the inequality of the mean shows that

\[
\mathrm{N} _ {p} (\mathbf {f}) \leqslant \left(\mathrm{N} _ {s} (\mathbf {f})\right) ^ {s / p} \left(\mathrm{N} _ {\infty} (\mathbf {f})\right) ^ {(p - s) / p}.\tag{12}
\]

Letting \(p\) tend to \(+\infty\) , it follows that

\[
\limsup _ {p \to + \infty} N _ {p} (\mathbf {f}) \leqslant N _ {\infty} (\mathbf {f})  .\tag{13}
\]

This proves that if \(+\infty \in I\) then \(J\) contains arbitrarily large numbers; thus \(I\) is indeed an interval of \(\overline{\mathbf{R}}\) , and \(\overline{I} = \overline{J}\) . The proposition will be proved if we show that \(p \mapsto N sub p(\mathbf{f})\) is continuous on \(\overline{J}\) , and it suffices to establish continuity at the end-points of \(J\) . We can suppose, moreover, that \(J\) does not reduce to a point. Let \(r\) and \(s\) be the left and right end-points of \(J\) ( \(r < s \leqslant +\infty\) ). Let \(A\) be the (measurable) set of \(x \in X\) such that \(|\mathbf{f}(x)| \geqslant 1\) ; then

\[
\int | \mathbf {f} | ^ {p} d | \mu | = \int | \mathbf {f} | ^ {p} \varphi_ {\mathrm{A}} d | \mu | + \int | \mathbf {f} | ^ {p} \varphi_ {\mathbf {C A}} d | \mu |.
\]

As \(p \in \mathbf{J}\) tends to \(r\) , \(|\mathbf{f}|^p\varphi_{\mathbf{A}}\) tends to \(|\mathbf{f}|^r\varphi_{\mathbf{A}}\) while decreasing, and \(|\mathbf{f}|^p\varphi_{\mathbf{C}_{\mathbf{A}}}\) tends to \(|\mathbf{f}|^r\varphi_{\mathbf{C}_{\mathbf{A}}}\) while increasing. Thus \(\int |\mathbf{f}|^p\varphi_{\mathbf{C}_{\mathbf{A}}}d|\mu|\) tends to \(\int^{*}|\mathbf{f}|^{r}\varphi_{\mathbf{C}_{\mathbf{A}}}d|\mu|\) (§1, No.~3, Th. 3). On the other hand, \(|\mathbf{f}|^p\varphi_{\mathbf{A}}\) is integrable for \(p \in \mathbf{J}\) , and \(\int |\mathbf{f}|^p\varphi_{\mathbf{A}}d|\mu|\) tends to \(\int |\mathbf{f}|^r\varphi_{\mathbf{A}}d|\mu|\) (§4, No.~3, Prop. 4). Therefore \(\int |\mathbf{f}|^p d|\mu|\) tends to \(\int^{*}|\mathbf{f}|^{r}d|\mu|\) , which proves the continuity of \(p \mapsto \mathrm{N}_p(\mathbf{f})\) at \(r\) .

The same reasoning may be applied at the point \(s\) if \(s < +\infty\) . Finally, suppose that \(s = +\infty\) . In view of (13), it suffices to prove that

\[
\liminf _ {p \to + \infty} N _ {p} (\mathbf {f}) \geqslant N _ {\infty} (\mathbf {f}).
\]

Now, let \(a\) be a number such that \(0 < a < \mathrm{N}_{\infty}(\mathbf{f})\) . Since, by hypothesis, there exist finite values of \(p\) such that \(\mathrm{N}_p(\mathbf{f}) < +\infty\) , the set \(A\) of \(x \in X\) such that \(|\mathbf{f}(x)| \geqslant a\) , which is measurable and non-negligible, is integrable by virtue of the inequality \(\varphi_A \leqslant (|\mathbf{f}|/a)^p\) ; moreover, we infer from this inequality that \(\mathrm{N}_p(\mathbf{f}) \geqslant a \cdot (|\mu|(\mathrm{A}))^{1/p}\) ; letting \(p\) tend to \(+\infty\) , it follows that \(\liminf_{p \to +\infty} \mathrm{N}_p(\mathbf{f}) \geqslant a\) , which completes the proof.

COROLLARY. --- If r, s, p are three numbers such that

\[
1 \leqslant r <   p <   s \leqslant + \infty ,
\]

then the intersection \(\mathcal{L}_{\mathrm{F}}^{r} \cap \mathcal{L}_{\mathrm{F}}^{s}\) is contained in \(\mathcal{L}_{\mathrm{F}}^{p}\) .

Note that in general the topologies induced on the intersection \(\mathcal{L}_{\mathrm{F}}^{r} \cap \mathcal{L}_{\mathrm{F}}^{s}\) by the topologies of the \(\mathcal{L}_{\mathrm{F}}^{p}\) ( \(r < p < s\) ) are distinct. If no further hypothesis is made on \(\mu\) , the topologies induced on \(\mathcal{L}_{\mathrm{F}}^{r} \cap \mathcal{L}_{\mathrm{F}}^{s}\) by those of \(\mathcal{L}_{\mathrm{F}}^{r}\) and \(\mathcal{L}_{\mathrm{F}}^{s}\) are in general not comparable (in other words, the ratio \(\mathrm{N}_{r}(\mathbf{f})/\mathrm{N}_{s}(\mathbf{f})\) can take arbitrarily large values and arbitrarily small values in \(L_{F}^{r} \cap L_{F}^{s}\) ; cf.~Exer. 8).

Prop. 4 may be sharpened when \(\mu\) is a bounded measure:

PROPOSITION 5. --- Let \(\mu\) be a bounded measure, and let \(\mathbf{f}\) be a \(\mu\) -measurable function with values in a Banach space \(\mathbf{F}\) . The set I of numbers \(p\) such that \(1 \leqslant p \leqslant +\infty\) and \(\mathrm{N}_p(\mathbf{f})\) is finite, is either empty or is an interval with left end-point \(p = 1\) and containing this point; moreover, \((|\mu|(\mathrm{X}))^{-1/p} \mathrm{N}_p(\mathbf{f})\) is an increasing function of \(p\) on I.

This is an immediate consequence of Prop. 4 above and of the Cor. of Prop. 4 of Ch. I, No.~3.

COROLLARY. --- If the measure \(\mu\) is bounded, the relation r < s implies \(L_{F}^{s} \subset L_{F}^{r}\) ; moreover, the topology of convergence in mean of order s is finer than the topology of convergence in mean of order r (on \(L_{F}^{s}\) ).

One can show that in general the topology of convergence in mean of order \(s\) is strictly finer than the topology of convergence in mean of order \(r\) (Exer. 8).

PROPOSITION 6. --- Let X be a discrete space, \(\mu\) the measure on X defined by placing a mass \(+1\) at each point of X. If \(\mathbf{f}\) is a mapping of X into the Banach space F, the set I of numbers \(p\) such that \(1 \leqslant p \leqslant +\infty\) and \(\mathrm{N}_p(\mathbf{f})\) is finite, is either empty or is an interval with right end-point \(+\infty\) and containing this point; moreover, \(\mathrm{N}_p(\mathbf{f})\) is a decreasing function of \(p\) on I.

For, \(\mu^{*}(|\mathbf{f}|) = \sum_{x\in X}|\mathbf{f}(x)|\) for every function \(\mathbf{f}\) ( \(\S 1\) , No.~1, Example), and \(\mathrm{N}_{\infty}(\mathbf{f}) = \| \mathbf{f}\| = \sup_{x\in X}|\mathbf{f}(x)|\) ; if there exists a number \(\alpha >0\) such that \(|\mathbf{f}(x)|\geqslant \alpha\) for infinitely many values of \(x\in X\) , then \(\mathrm{N}_p(\mathbf{f}) = +\infty\) for every finite \(p\) ; in the contrary case, there exists an \(x_0\in X\) such that \(|\mathbf{f}(x_0)| = \| \mathbf{f}\|\) , whence

\[
\mathrm{N} _ {\infty} (\mathbf {f}) = | \mathbf {f} (x _ {0}) | \leqslant \mathrm{N} _ {p} (\mathbf {f})
\]

for every finite p.~Since the function \(\log\mathrm{N}_{p}(\mathbf{f})\) is convex with respect to 1/p and takes its smallest value at the point \(+\infty\) , it is necessarily a decreasing function of p on I (FRV, I, §4, No.~3, Prop. 5), which completes the proof.

COROLLARY. --- If X is discrete and the measure \(\mu\) is defined by a mass +1 at each point of X, then the relation r < s implies \(L_{F}^{r} \subset L_{F}^{s}\) ; moreover, the topology of convergence in mean of order r is finer than the topology of convergence in mean of order s (on \(L_{F}^{r}\) ).

